\chapter{Data Architecture}
\label{ch:data}

% Sources: thesis-pack/06-data-layer/DATABASE_IMPLEMENTATION.md (sec.3 schema,
% sec.4 phases 3-5, sec.6 guardrails), adr/0003, adr/0005,
% 04-architecture-agentic/agentic-core-README.md sec.13-14,
% 10-instrument/evaluation-README.md. Written as built: PostgreSQL only.
% Identifiers: table names and event kinds only, no code names.

Chapter~\ref{ch:agentic} showed that the ten agents never talk to each other
directly and that every message, state change and decision passes through one
PostgreSQL database (Figure~\ref{fig:ag:architecture}). In this chapter we
describe that database. We show how it stores the actors, their states, the
events between them and the resources they use. We also show how the same
structures serve two purposes at once: they run the simulation, and they are
the record that the evaluation scores.

% =====================================================================
\section{Design Principles}
\label{sec:data:principles}
% Source: ADR-0003 (Decision, Explicitly NOT claimed), ADR-0005 (Decision,
% Consequences).

The database has three jobs in the system.
\begin{enumerate}
  \item It is the \textbf{state store}. The current state of every agent lives
  in the database and nowhere else. An agent that restarts reads its state back
  from there.
  \item It is the \textbf{message broker}. A message from one agent to another
  is a row in a table. The database tells the recipient that the row has
  arrived.
  \item It is the \textbf{scoring backbone}. Every agent, and the workflow
  baseline as well, writes what it received and what it decided into one
  append-only log. The evaluation computes every score from that log and from
  nothing else.
\end{enumerate}
This is what we mean by \emph{database-centric}. The database is the single
source of truth, and all three jobs read and write the same store. There is no
second copy of any state in a file or in memory that could drift away from it.

We keep one limit clear from the start. Using a shared store to coordinate
independent programs is an old idea. Blackboard systems \parencite{nii1986}
and tuple spaces \parencite{gelernter1985} did so decades before this work
(Section~\ref{sec:bg:dbcentric}).
We therefore do not claim a shared database as a new way for agents to
coordinate. What we do claim concerns evaluation. One append-only log is shared
by both systems under test and by the scoring code. This makes the comparison
fair, checkable and repeatable. Section~\ref{sec:data:evalbackbone} makes this
argument in full. The log follows the idea of event sourcing, where the history
of events is stored and every later view is derived from it
\parencite{fowler2005,overeem2021}.

A single source of truth has a clear cost. If the database is down, the whole
simulation is down, because no agent can load its state or receive a message.
We accept this. Before the runtime starts the ten agent programs, it checks
that the database answers, and it refuses to start if it does not.

% =====================================================================
\section{Data Model}
\label{sec:data:conceptual}
% Source: DATABASE_IMPLEMENTATION.md sec.3 (12 tables as built).

The model has four kinds of thing: actors, their states, the events between
them, and the physical resources they use. The ten agents are the actors.
Vehicles are not actors. A bus, the patrol car and the ambulance are rows in a
fleet table that a person drives (Section~\ref{sec:ag:overview}). Scenarios are
not stored in the database. They are files that a dispatcher reads and turns
into messages at the right simulated time.

The database holds twelve tables. We group them by the job they do
(Table~\ref{tab:data:tables} and Figure~\ref{fig:data:tables}).

\begin{description}
  \item[Runtime tables.] Five tables keep the simulation running: the agent
  states, the mailboxes, the world reports, the prompt records and the timers.
  They are cleared between two runs.
  \item[The audit table.] The event log is the only table that is never
  changed or cleared. It holds the history of every run.
  \item[Sensing tables.] Six small tables hold the facts that some agents
  perceive rather than hear about in a message. Examples are an active flood
  alert, the remaining beds of a shelter, a road closure and the status of each
  vehicle. The sensing step of BIS, Hochbahn, the police officer and the
  paramedic reads them.
\end{description}

% The twelve tables in three groups, and who writes and reads each group.
% Source: thesis-pack/06-data-layer/DATABASE_IMPLEMENTATION.md sec.3.
% Input from inc/data_architecture.tex, Section "Data Model".
\begin{figure}[htbp]
\centering
\begin{tikzpicture}[
  tbl/.style={draw, rounded corners=2pt, fill=orange!8, minimum width=3.1cm,
    minimum height=0.5cm, align=center, font=\footnotesize\ttfamily},
  audit/.style={draw, thick, rounded corners=2pt, fill=orange!20, minimum width=3.1cm,
    minimum height=0.6cm, align=center, font=\footnotesize\ttfamily},
  head/.style={align=center, font=\footnotesize\bfseries, minimum height=0.8cm},
  actor/.style={draw, rounded corners=2pt, fill=blue!6, text width=3.1cm,
    minimum height=0.8cm, align=center, font=\footnotesize},
  group/.style={draw, dashed, rounded corners=3pt, inner sep=4pt},
  wr/.style={-{Stealth[length=1.8mm]}, thick, gray!70!black},
  lbl/.style={font=\scriptsize, align=center, text=gray!30!black},
]
  % --- runtime column ---------------------------------------------------
  \node[head] (hrt) at (0,0) {Runtime\\[-1pt]{\scriptsize\normalfont cleared before each run}};
  \node[tbl, below=0.1cm of hrt] (states) {states};
  \node[tbl, below=0.12cm of states] (msg) {messages};
  \node[tbl, below=0.12cm of msg] (world) {world\_events};
  \node[tbl, below=0.12cm of world] (side) {tick\_sidecars};
  \node[tbl, below=0.12cm of side] (cron) {cronjobs};
  \node[group, fit=(hrt)(cron)] (rt) {};

  % --- audit column -----------------------------------------------------
  \node[head] (hau) at (4.4,0) {Audit\\[-1pt]{\scriptsize\normalfont never cleared}};
  \node[audit, below=0.1cm of hau] (log) {event\_log};
  \node[lbl, below=0.1cm of log] (loglbl) {append-only,\\enforced by a trigger;\\every row names its\\run, agent and model};
  \node[group, fit=(hau)(loglbl)] (au) {};

  % --- sensing column ---------------------------------------------------
  \node[head] (hse) at (8.8,0) {Sensing\\[-1pt]{\scriptsize\normalfont read by the agents' sensing step}};
  \node[tbl, below=0.1cm of hse] (flood) {flood\_alert};
  \node[tbl, below=0.12cm of flood] (weather) {weather\_alert};
  \node[tbl, below=0.12cm of weather] (shelter) {shelter\_capacity};
  \node[tbl, below=0.12cm of shelter] (fleet) {fleet};
  \node[tbl, below=0.12cm of fleet] (road) {road\_closure};
  \node[tbl, below=0.12cm of road] (med) {medical\_incident};
  \node[group, fit=(hse)(med)] (se) {};

  % --- writers, top -----------------------------------------------------
  \node[actor, above=1.0cm of rt.north] (agents) {Ten agents};
  \node[actor, above=1.0cm of au.north] (base) {Workflow baseline};
  \draw[wr, <->] (agents.south) -- node[lbl, left] {read, write} (rt.north);
  \draw[wr] (agents.south east) -- node[lbl, above, sloped, pos=0.45] {write} (au.north west);
  \draw[wr] (base.south) -- node[lbl, right] {write} (au.north);
  \draw[wr] (se.north) |- ([yshift=0.45cm]agents.north)
    node[lbl, above, pos=0.75] {read sensing tables} -- (agents.north);

  % --- other parties, bottom -------------------------------------------
  \node[actor, below=1.0cm of rt.south] (ui) {Interface and\\world client};
  \node[actor, below=1.0cm of au.south -| log, xshift=-0.2cm] (harness) {Evaluation harness};
  \node[actor, below=1.0cm of se.south] (score) {Scorer, rubric,\\no-replan check};
  \draw[wr, <->] (ui.north) -- node[lbl, left] {watch states,\\write world reports} (rt.south);
  \draw[wr] (harness.north) -- node[lbl, right] {manifest} (harness.north |- au.south);
  \draw[wr] (au.south east) -- node[lbl, above, sloped] {read} (score.north west);
\end{tikzpicture}
\caption[The twelve tables of the database]{The twelve tables of the database, grouped by job. The runtime tables
are cleared before each run. The event log is append-only and holds every run.
Both systems write to it, and the scoring code reads only from it.}
\label{fig:data:tables}
\end{figure}


\subsection{Tables}
\label{sec:data:schema}

\begin{table}[htbp]
\centering
\caption{The twelve tables, grouped by job.}
\label{tab:data:tables}
\small
\begin{tabularx}{\textwidth}{>{\raggedright\arraybackslash}p{3.7cm}>{\raggedright\arraybackslash}X>{\raggedright\arraybackslash}p{3.0cm}}
\toprule
Table & What one row holds & Life cycle \\
\midrule
\multicolumn{3}{l}{\emph{Runtime}} \\
\texttt{states} & The full state of one agent, as one document & Overwritten on every save \\
\texttt{messages} & One message: sender, recipient, type, body, and the times it arrived and was read & Marked as read, never deleted during a run \\
\texttt{world\_events} & One report from the world that an action has finished or a vehicle has stopped & Marked as read \\
\texttt{tick\_sidecars} & The full prompt and the full answer of one model call & Written before and after the call \\
\texttt{cronjobs} & One timer an agent has set & Record only \\
\midrule
\multicolumn{3}{l}{\emph{Audit}} \\
\texttt{event\_log} & One thing an agent received or decided & Append-only, never cleared \\
\midrule
\multicolumn{3}{l}{\emph{Sensing}} \\
\texttt{flood\_alert} & An alert for a facility, with the number of people to evacuate & Read by sensing \\
\texttt{weather\_alert} & A weather warning & Read by sensing \\
\texttt{shelter\_capacity} & The remaining beds of one shelter & Read by sensing \\
\texttt{fleet} & One vehicle: its driver, type, status, mission, position and arrival time & Read by sensing \\
\texttt{road\_closure} & A closed road & Read by sensing \\
\texttt{medical\_incident} & A medical incident and whether it is resolved & Read by sensing \\
\bottomrule
\end{tabularx}
\end{table}

\paragraph{One document per agent.} The \texttt{states} table has one row per
agent. The row holds the agent's whole state as one JSON document: its inbox,
its current step, its remaining plan, its history of finished steps, its view
of the fleet and its outgoing messages. Every save replaces the whole row. The
row also carries the run identifier and the model name, so that a state can
always be tied to the run and the model that produced it.

We did not split this document into separate tables. The runtime always loads
and saves an agent's state as a whole, so it gains nothing from a finer
structure. The evaluation does not read this document either. It reads the
event log (Section~\ref{sec:data:eventlog}), which records every change to the
step and the plan as it happens. A normalised state would add joins and
migration work without serving a single measurement.

\paragraph{One fleet table for three kinds of vehicle.} The \texttt{fleet}
table carries both a bus number and a vehicle number. Hochbahn's sensing looks
up buses by the first, and the police officer and the paramedic look up their
cars by the second. One table serves all three readers.

% =====================================================================
\section{Append-Only Event Log}
\label{sec:data:eventlog}
% Source: agentic-core-README.md sec.14 (envelope, kinds, KPI fields);
% DATABASE_IMPLEMENTATION.md sec.3 (unique key, trigger, pk ordering) and
% Phase 3 (reset excludes event_log); evaluation-README.md (manifest rows).

The event log is the record of what happened in a run. Every agent writes to
it, and no one may change what has been written.

\paragraph{The row.} Each row has the same seven fields:
\begin{itemize}
  \item an \textbf{identifier} that counts up within one agent's log
  (\texttt{evt-0}, \texttt{evt-1}, and so on);
  \item the \textbf{wall-clock time} and the \textbf{simulated time} of the
  entry;
  \item a \textbf{kind} and a \textbf{subkind}, which say what sort of entry it
  is;
  \item a one-line \textbf{summary} for a human reader;
  \item a \textbf{details} document with everything else.
\end{itemize}
Three more columns place the row: the \textbf{run}, the \textbf{agent} and the
\textbf{model}. With them, one table holds every run of every agent, and any
single run or agent can be selected from it.

\paragraph{The kinds of entry.} Five kinds cover everything the systems write.
\begin{description}
  \item[\texttt{event\_received}] An event woke the agent: a message, a
  finished step or a timing conflict. The details hold the event as it arrived.
  \item[\texttt{llm\_response}] The agent made a decision with its model. The
  details hold the reasoning, the committed step, the remaining plan, the
  outgoing messages, the token counts and the latency. A dispatch by the clock
  that needs no model uses the same kind, with the subkind \texttt{cron\_fire}
  and zero tokens.
  \item[\texttt{world\_ack}] The world confirmed that an action had finished.
  The details hold the vehicle and its real position.
  \item[\texttt{decision}] The workflow baseline made a decision. This is the
  baseline's counterpart to \texttt{llm\_response}
  (Section~\ref{sec:data:evalbackbone}).
  \item[\texttt{run\_manifest}, \texttt{run\_outcome}] The evaluation harness
  recorded how a run was set up and how it ended
  (Section~\ref{sec:data:replay}).
\end{description}

\paragraph{Measurement fields.} Every decision row carries the same five
fields for the evaluation, whether or not anything went wrong. They are the
list of rule violations the answer tried, their count, a fingerprint of what
the model saw, the retry count and a failure code. Because the fields are
always present, a zero means that we measured zero. It does not mean that the
value is missing. The fingerprint is a hash over the agent's observable
situation, taken before the model call. It lets the evaluation see whether two
decisions were made from the same observation (Section~\ref{sec:ag:engine}).

\paragraph{Three guarantees.} The database, not a convention in the code,
enforces how the log behaves.
\begin{enumerate}
  \item \textbf{Append-only.} A trigger rejects every update and every delete
  on the table. Once a row is written, it stays as it is. A run cannot be
  quietly edited or removed after the fact.
  \item \textbf{Written once.} The run, the agent and the row identifier
  together form a unique key. An agent saves its whole in-memory log each time
  it saves its state, and rows that are already stored are skipped. Saving
  twice therefore never creates a duplicate.
  \item \textbf{Read in order.} Each row also gets a sequence number when it is
  inserted, and the log is read back in that order. Sorting by the identifier
  would put \texttt{evt-10} before \texttt{evt-2}.
\end{enumerate}

Between two runs, the runtime clears the runtime tables so that no agent
starts with the state or the unread mail of the run before. The event log is
left out of this reset on purpose. It is divided by run, so old runs do not
disturb a new one, and the trigger would reject the delete in any case.

% =====================================================================
\section{The Event Log as Evaluation Backbone}
\label{sec:data:evalbackbone}
% Source: ADR-0003 (Decision), agentic-core-README.md sec.14 (decision kind),
% evaluation-README.md sec.5b (harness invisible to scorer, facts.py test).
% This is contribution 4 in sec:intro:contributions.

The event log is shared by three parties: the agentic system, the workflow
baseline and the code that scores them. Figure~\ref{fig:data:backbone} shows
this arrangement. It is the methodological contribution of the data layer.

% The event log as the evaluation backbone: both systems write the same rows,
% one scorer reads them, harness rows are nested so they cannot change a score.
% Source: thesis-pack/adr/0003, agentic-core-README.md sec.14,
% 10-instrument/evaluation-README.md. Input from inc/data_architecture.tex,
% Section "The Event Log as Evaluation Backbone".
\begin{figure}[htbp]
\centering
\begin{tikzpicture}[
  sys/.style={draw, rounded corners=2pt, fill=blue!6, text width=3.0cm,
    minimum height=1.1cm, align=center, font=\footnotesize},
  base/.style={draw, rounded corners=2pt, fill=gray!10, text width=3.0cm,
    minimum height=1.1cm, align=center, font=\footnotesize},
  db/.style={draw, thick, rounded corners=3pt, fill=orange!8, text width=4.1cm,
    align=left, font=\footnotesize, inner sep=6pt},
  read/.style={draw, rounded corners=2pt, fill=green!6, text width=2.7cm,
    minimum height=0.62cm, align=center, font=\footnotesize},
  wr/.style={-{Stealth[length=1.8mm]}, thick, gray!70!black},
  lbl/.style={font=\scriptsize, align=center, text=gray!30!black},
]
  % --- the log in the centre ------------------------------------------
  \node[db] (log) {%
    \textbf{\texttt{event\_log}}\\[1pt]
    \emph{append-only, one table}\\[3pt]
    same seven fields per row\\
    plus run, agent, model\\[3pt]
    \texttt{event\_received}\\
    \texttt{llm\_response}\\
    \texttt{world\_ack}\\
    \texttt{decision}\\
    \texttt{run\_manifest}, \texttt{run\_outcome}};

  % --- writers, left ----------------------------------------------------
  \node[sys, left=1.1cm of log.north west, anchor=north east] (ag) {%
    \textbf{Agentic system}\\ten agents, shared engine\\[2pt]
    {\scriptsize writes \texttt{llm\_response}}};
  \node[base, left=1.1cm of log.south west, anchor=south east] (wf) {%
    \textbf{Workflow baseline}\\one incident commander\\[2pt]
    {\scriptsize writes \texttt{decision}}};
  \draw[wr] (ag.east) -- (log.west |- ag.east);
  \draw[wr] (wf.east) -- (log.west |- wf.east);

  % --- harness, top -----------------------------------------------------
  \node[base, above=0.9cm of log, text width=3.6cm, minimum height=0.8cm] (h) {%
    \textbf{Evaluation harness}};
  \draw[wr] (h.south) -- node[lbl, right] {manifest, nested one\\level down in details} (log.north);

  % --- readers, right ---------------------------------------------------
  \node[read, right=1.1cm of log.north east, anchor=north west] (sc) {Deterministic scorer};
  \node[read, below=0.35cm of sc] (ru) {Held-out rubric};
  \node[read, below=0.35cm of ru] (g0) {No-replan check};
  \node[lbl, below=0.2cm of g0] (one) {one reader function,\\blind to the system};
  \foreach \r in {sc, ru, g0} {
    \draw[wr] (log.east |- \r.west) -- (\r.west);
  }

  % --- guarantees, bottom ------------------------------------------------
  \node[lbl, below=0.35cm of log, text width=6.5cm] {a trigger rejects every update and
    delete; a unique key on run, agent and row makes a second save harmless};
\end{tikzpicture}
\caption[The event log as the backbone of the evaluation]{The event log as the backbone of the evaluation. Both systems write
rows of the same format into one append-only table. The baseline's decision
rows carry their own kind because no model was called. The harness records
each run's setup in the same log, one level deeper than any field the scorer
reads. The scorer, the rubric and the no-replan check read the same rows for
both systems.}
\label{fig:data:backbone}
\end{figure}


\paragraph{Both systems write the same rows.} The agentic agents write their
rows through the shared engine. The workflow baseline has no agent state and
no model, so it uses a small writer of its own. That writer produces the same
row format, with the same numbering and the same fields, into the same table.
There is one difference, and we made it on purpose. The baseline's decision
rows have the kind \texttt{decision} and an empty model name, and their token
counts and latency are zero. We did not label them \texttt{llm\_response},
because no model was called, and that label would misreport the control
condition. The baseline computes its observation fingerprint by the same rule
as the agents.

\paragraph{One reader for both.} The scorer reads the log through one
function, and that function runs in the same way for both systems. It does not
know which system wrote a row. It sees only rows, their kinds and their
details. The rubric and the no-replan check (Section~\ref{sec:eval:noreplan})
read the same rows.

\paragraph{The harness cannot change a score.} The harness writes its own rows
into the log, to record how each run was set up. Their content sits one level
deeper inside the details than the fields the scorer reads. A manifest that
carried a route length at the top level, for example, would add kilometres to
the run it describes. A test checks that scoring a run with its manifest rows
gives the same result as scoring it without them.

\paragraph{What the shared log gives the evaluation.}
\begin{itemize}
  \item \textbf{Deterministic scoring.} Every outcome is a fixed function of
  the log (Section~\ref{sec:eval:scorer}). No human judges a run, and no model
  judges one either. The same log always gets the same score.
  \item \textbf{Checkable parity.} Because both systems write the same rows
  into the same table, the claim that they are comparable can be tested
  against shared code instead of stated. An executable parity test suite
  does this, and lists the seven places where the two systems still differ.
  \item \textbf{Re-scoring.} The log keeps every run. When a scoring rule
  changes, every run can be scored again from its rows, without running it
  again. Because the log is append-only, a run that is dropped from an
  analysis still leaves its rows behind.
\end{itemize}

We repeat the limit from Section~\ref{sec:data:principles}. The contribution is
the log as the shared backbone of a comparison and its scoring. The shared
database as a coordination medium is not new, and we do not present it as new.

% =====================================================================
\section{Runtime State and Message Passing}
\label{sec:data:runtime}
% Source: DATABASE_IMPLEMENTATION.md Phase 2-3 (drain marks consumed_at, FOR
% UPDATE SKIP LOCKED, alias folding, dead letters, reset), ADR-0005 (audit
% namespace, no run_id on messages/world_events), agentic-core-README.md sec.13
% (sidecars twice per wake).

\paragraph{Saving state.} An agent saves its state after every wake. In one
transaction, the save replaces the agent's row in \texttt{states} and appends
any new rows of its log to \texttt{event\_log}. After the transaction has
committed, the agent announces the change (Section~\ref{sec:data:notify}).

\paragraph{Sending and reading messages.} To send a message, an agent inserts a
row with the recipient's key into the \texttt{messages} table. To read its mail, the
recipient selects its unread rows and marks them as read. It does not delete
them. The mailbox therefore keeps a full record of every message and of when it
was read. The rows are locked while they are taken, and rows locked by another
reader are skipped, so no message is ever read twice.\footnote{This is
PostgreSQL's \texttt{SELECT \dots\ FOR UPDATE SKIP LOCKED}, see
\url{https://www.postgresql.org/docs/current/sql-select.html}.}

\paragraph{One address per agent.} Some senders name a recipient by a short
alias, such as a vehicle call sign, rather than by the recipient's key. All
messages pass through one insert step, and that step maps every alias onto the
recipient's key. If a recipient cannot be resolved, the message is logged as
undeliverable instead of being lost. Since every message passes this one
point, a test follows every alias from sender to recipient.

\paragraph{Messages that are recorded but not delivered.} Each sender has a
routing table. For some recipients it marks a message as \emph{audit}, meaning
that it should be kept for the record but not delivered. Such a message is
stored under a separate key space that no agent reads. It stays available to
the evaluation and never wakes an agent.

\paragraph{World reports.} When a vehicle arrives, or stops half-way because
its plan changed, the world writes a row to \texttt{world\_events}. In a live
run the source is the map in the browser. In the evaluation it is the headless
world client (Section~\ref{sec:eval:harness}). The agent reads these rows in
the same way as its mail and takes its real position from them
(Section~\ref{sec:ag:runtime}).

\paragraph{Prompt records and timers.} Each model call writes one row to
\texttt{tick\_sidecars} before the call and updates it after. The first write
lets the interface show that a call is in progress. The second adds the
model's full answer. Token counts and latency are not stored here. They are in
the decision row of the event log. The \texttt{cronjobs} table records the
timers that an agent sets, but it does not run them. The clock that fires the
timers stays inside each agent program. Moving it into the database would
change how and when an agent acts, and that behaviour is part of what we
measure.

\paragraph{Separation between runs.} The \texttt{messages} and
\texttt{world\_events} tables have no run column. What keeps one run's mail
out of the next run is the reset that clears these tables before every run.
The harness performs this reset at the start of each cell. A run started
without it would inherit unread mail from the run before. We record this as a
limitation of the schema. The event log is not affected, because it carries
the run on every row.

% =====================================================================
\section{Live Notifications}
\label{sec:data:notify}
% Source: DATABASE_IMPLEMENTATION.md Phase 5 (three channels, application-side
% NOTIFY after commit, NOTIFY_FALLBACK_POLL_S = 5 s default, SelectorEventLoop
% thread on Windows, error backoff unchanged).

An agent should wake as soon as mail arrives, without asking the database
again and again. Reactive tuple spaces solve the same problem by running code
when new data arrives \parencite{omicini1999,cabri1998}. PostgreSQL offers a
publish-and-subscribe mechanism for this, called LISTEN and
NOTIFY.\footnote{\url{https://www.postgresql.org/docs/current/sql-notify.html}}
A program listens on a named channel, and any other program can send a short
notification on that channel.

Each agent has three channels: one for its state, one for its mailbox and one
for its world reports. Whoever writes a message, a world report or a state
sends a notification on the matching channel. The notification is sent from
the writing program after its transaction has been committed, not from a
trigger inside the database. By the time the recipient wakes and reads, the
new row is therefore visible to it. Every writer goes through the same insert
functions, so no writer can forget to notify.

A notification can still be lost, for example when a connection drops. We
therefore kept a slow poll as a safety net. An agent waits for a notification
or for five seconds, whichever comes first, and then looks at its mailbox.
A lost notification makes an agent slower, but it never leaves the agent
stuck. No part of the system depends on a notification arriving to be correct.

The listener runs on its own thread with its own event loop. The database
driver cannot listen asynchronously on the event loop that Python uses by
default on Windows, and the agents need that default loop for other reasons.
The same code path runs on every operating system.

The observation interface uses the same mechanism. It listens on the state
channel of every agent and redraws the map when a state changes
(Chapter~\ref{ch:ui}).

% =====================================================================
\section{External Data Sources}
\label{sec:data:external}
% Source: ADR-0005 decision 5 (working directories keep the caches),
% evaluation-README.md (route_km recorded at dispatch on both sides; route
% cache is one file), agentic_system.tex sec:ag:crons (fallback).

Two kinds of data come from outside the database. Addresses are turned into
coordinates by a geocoder, and drives are routed over real Hamburg roads with
OpenRouteService. The places themselves, meaning the Augustinum, the two
schools and the depots, are part of the case (Chapter~\ref{ch:casestudy}).

Both services are called through caches on disk. Each agent keeps these caches
in its own working directory, together with its program log. No state is kept
there. The caches serve two ends. They keep the number of calls to the routing
service low, and a repeated query returns the stored route. If
routing fails, the runtime falls back to a straight-line estimate and marks
the action as degraded, as Section~\ref{sec:ag:crons} describes.

The route length also enters the evaluation. When a bus is dispatched, both
systems record the kilometres of its route, measured along the same routed
path. The scorer reads this value from the log for both sides in the same way.

% =====================================================================
\section{Determinism and Replay}
\label{sec:data:replay}
% Source: evaluation-README.md (RESCUESIM_RUN_ID pin, manifest rows, seeded not
% reproducible, :cloud tags, serial sweep reasons), ADR-0014 (scale 8.4 only).

A run can be scored again at any time if its log is complete. Running it
again gives a similar run, but not the same one. We record the following so
that re-scoring is always possible.

\paragraph{One run identifier.} All ten agent programs of one run share one run
identifier. The runtime sets it when a run starts and passes it to every agent.
Every row of the event log carries it, so the ten agents' rows of one run can
be joined into a single history.

\paragraph{The manifest inside the log.} At the start of each run the harness
writes a manifest row into the event log. It names the system, the model and
its digest, the scenario and its hash, the time scale, the seed, the code
version and the hash of the scoring thresholds. Section~\ref{sec:eval:harness}
lists the full content. Because the manifest sits in the append-only log,
the setup of a run cannot be separated from the run.

\paragraph{Seeded, but not reproducible.} Every model call receives a seed that
is derived from the scenario, the time scale and the replicate number. The seed
is logged. A seeded call to a language model is still only close to
deterministic, because batching and cache state inside the inference server
change the sampling \parencite{yuan2025}. We therefore claim that runs are seeded and logged, not
that they can be reproduced at the level of each model answer. This is also why
every cell is repeated (Section~\ref{sec:eval:analysis}). Models served from a
hosted endpoint cannot be pinned to one version. For them, the model's details
are captured with every run, so that a change of model would at least be
visible.

\paragraph{The time scale.} The ratio of simulated to wall-clock time is a
logged constant of each run. Every scored run used the value 8.4. The log
records it so that a run could be priced again at another scale, but this
thesis makes no claim about other scales (Section~\ref{sec:eval:responsetime}).

\paragraph{Why runs cannot run in parallel.} The data layer sets one practical
limit on the evaluation. The \texttt{states} table has one row per agent key,
and the agent keys are fixed names. Two runs on one database would overwrite
each other's state. The agents' working directories are fixed paths as well,
and the route cache is a single file. The evaluation therefore runs its cells
one after another (Section~\ref{sec:eval:harness}).

% =====================================================================
\section{Chapter Summary}
\label{sec:data:summary}

One PostgreSQL database is the state store, the message broker and the scoring
backbone of the whole system. Twelve tables serve these jobs. Five runtime
tables keep agent states, mailboxes, world reports, prompt records and timers,
and they are cleared between runs. Six sensing tables hold what agents
perceive. The event log is append-only, enforced by a trigger, and every row
names its run, agent and model. Both the agentic system and the workflow
baseline write the same rows into it, and one scorer reads them for both.
Notifications wake agents without polling, and a five-second fallback keeps a
lost notification from stalling anything. We do not claim the shared database
as a new coordination medium. We claim the shared log as the backbone of a
fair, checkable and repeatable comparison. Chapter~\ref{ch:ui} describes the
interface that watches this database during a live run.
